% Para insertar imágenes
\usepackage{graphicx}

% Para tener referencias
\usepackage{hyperref}

% Para poner matemáticas
\usepackage{amsthm, amssymb}

% Comentarios en pdf
\setlength {\marginparwidth }{2cm}
\usepackage{todonotes}

% Matrices
\usepackage{amsmath}

% Para poner definiciones con color
\usepackage[most]{tcolorbox}

% Colores
\usepackage{xcolor}

% Quita sangrías/Pone saltos en párrafos
%\usepackage{parskip}

% Español
\usepackage[spanish]{babel}

% imagenes svg
\usepackage{svg}

% para insertar imágenes en una posición estricta
\usepackage{float}

% Definición personalizada de tcolorbox
\newtcolorbox{miDefinicion}[2][]{
  colback=teal!5!white,
  colframe=teal!75!black,
  coltitle=white,
  fonttitle=\bfseries,
  title=Definición: #2,
  arc=2mm,
  boxrule=0.8pt,
  enhanced,
  drop shadow,
  #1
}


% Para poner call outs chidos
\usepackage{mdframed}
% Definimos el estilo del callout con línea a la izquierda
\newmdenv[
  topline=false,      % Quita la línea superior
  bottomline=false,   % Quita la línea inferior
  rightline=false,    % Quita la línea derecha
  leftline=true,      % Activa solo la línea izquierda
  linecolor=cyan,     % Color de la línea (puedes usar blue, red, etc.)
  linewidth=3pt,      % Grosor de la línea
  innerleftmargin=10pt, % Espacio entre la línea y el texto
  innerrightmargin=0pt,
  innertopmargin=0pt,
  innerbottommargin=0pt
]{miencuadre}

\usepackage{booktabs}

\usepackage{dirtree}

\usepackage{varwidth}

\usepackage{fancyvrb}

\usepackage{pmboxdraw}
\DeclareUnicodeCharacter{251C}{\textSFviii} % ├
\DeclareUnicodeCharacter{2500}{\textSFx}    % ─
\DeclareUnicodeCharacter{2514}{\textSFii}   % └
\DeclareUnicodeCharacter{2502}{\textSFxi}   % │